\documentclass{article}
\usepackage[utf8]{inputenc}
\usepackage{hyperref}
\usepackage{graphicx}

\title{Web Application Documentation}
\author{Project Team}
\date{\today}

\begin{document}

\maketitle

\section{Introduction}
This document provides comprehensive documentation for the Web Application project, a full-stack web application built using modern technologies and best practices.

\section{Project Overview}
The application is a web-based system that provides user authentication and authorization capabilities. It demonstrates the implementation of secure user management with features like signup, login, and protected routes.

\section{Architecture}
The project follows a microservices architecture with three main components:

\subsection{Frontend Service}
\begin{itemize}
    \item Built with React 19 and TypeScript
    \item Uses React Router for client-side routing
    \item Implements protected routes for authenticated users
    \item Communicates with backend via RESTful API using Axios
    \item Features responsive design for various screen sizes
\end{itemize}

\subsection{Backend Service}
\begin{itemize}
    \item Built with NestJS framework
    \item Implements RESTful API endpoints
    \item Uses JWT for authentication
    \item Includes Swagger/OpenAPI documentation
    \item Follows modular architecture with separate modules for auth and user management
\end{itemize}

\subsection{Database Service}
\begin{itemize}
    \item MongoDB for data persistence
    \item Stores user data and authentication information
    \item Mongoose ODM for database interactions
    \item Supports data validation and schema enforcement
\end{itemize}

\section{Technologies Used}
\subsection{Frontend}
\begin{itemize}
    \item React 
    \item TypeScript
    \item React Router
    \item Axios
    \item Create React App
\end{itemize}

\subsection{Backend}
\begin{itemize}
    \item NestJS
    \item TypeScript
    \item Passport JWT
    \item MongoDB with Mongoose
    \item Swagger/OpenAPI
\end{itemize}

\subsection{Database}
\begin{itemize}
    \item MongoDB 4.4+
    \item Mongoose ODM
    \item MongoDB Atlas (optional for cloud deployment)
\end{itemize}

\subsection{DevOps}
\begin{itemize}
    \item Docker
    \item Docker Compose
    \item Nginx (for frontend serving)
    \item Node.js 18+
\end{itemize}

\section{API Documentation}

\subsection{Authentication Endpoints}
\begin{itemize}
    \item POST /auth/signup
    \begin{itemize}
        \item Creates new user account
        \item Required fields: username, email, password
    \end{itemize}
    \item POST /auth/login
    \begin{itemize}
        \item Authenticates user and returns JWT token
        \item Required fields: email, password
    \end{itemize}
\end{itemize}

\subsection{User Endpoints}
\begin{itemize}
    \item GET /user/profile
    \begin{itemize}
        \item Returns authenticated user's profile
        \item Requires valid JWT token
    \end{itemize}
    \item PUT /user/profile
    \begin{itemize}
        \item Updates user profile information
        \item Requires valid JWT token
    \end{itemize}
\end{itemize}

\section{Security}
The application implements several security measures:
\begin{itemize}
    \item JWT-based authentication
    \item Password hashing using bcrypt
    \item Protected routes in both frontend and backend
    \item CORS protection
    \item Environment-based configuration
\end{itemize}

\section{Deployment}
The application can be deployed using either Docker Compose or traditional npm-based deployment. Docker Compose is recommended for production environments as it ensures consistency across different deployment environments.

\subsection{Docker Compose Deployment}
Includes:
\begin{itemize}
    \item Frontend container (Node.js 18+ for build, Nginx for serving)
    \item Backend container (Node.js 18+)
    \item MongoDB container
    \item Docker network for service communication
    \item Volume for MongoDB data persistence
\end{itemize}

\subsection{Traditional Deployment}
Requirements:
\begin{itemize}
    \item Node.js 18 or later
    \item npm 8 or later
    \item MongoDB 4.4 or later
\end{itemize}

\end{document} 